\section{常用公式、性质、方程式与规律总结}

\subsection{物质的量}

\[
n = \frac{N}{N_A},\qquad n = \frac{m}{M},\qquad n = \frac{V}{V_m},\qquad n = c\cdot V
\]

\[
N_A = 6.02\times 10^{23}\ \text{mol}^{-1},\qquad V_m = 22.4\ \text{L/mol}\ (\text{标准状况})
\]

\textbf{阿伏加德罗定律}：同温同压下，相同体积的任何气体含有相同数目的分子。

\[
\frac{V_1}{V_2} = \frac{n_1}{n_2},\qquad \frac{\rho_1}{\rho_2} = \frac{M_1}{M_2}
\]

\textbf{物质的量浓度与质量分数换算}：
\[
c = \frac{1000\rho\omega}{M}
\]
（$\rho$ 单位 $\text{g/mL}$，$\omega$ 为质量分数，$M$ 为摩尔质量）

\textbf{稀释定律}：$c_1V_1 = c_2V_2$

\textbf{易错}：
\begin{itemize}
  \item 使用 $22.4\ \text{L/mol}$ 时务必确认标准状况（$0^\circ\text{C}$，$101\ \text{kPa}$）且为气体
  \item 气体摩尔体积适用于气体（包括混合气体），但不适用于固体、液体
  \item 计算物质的量浓度时注意溶液体积以 $\text{L}$ 为单位
\end{itemize}

\subsection{氧化还原反应}

\textbf{核心规律}：
\begin{itemize}
  \item 升（化合价升高）$\to$ 失（失去电子）$\to$ 氧（被氧化）$\to$ 还（作还原剂）
  \item 降（化合价降低）$\to$ 得（得到电子）$\to$ 还（被还原）$\to$ 氧（作氧化剂）
  \item 氧化性：氧化剂 $>$ 氧化产物
  \item 还原性：还原剂 $>$ 还原产物
\end{itemize}

\textbf{守恒法配平}：
\begin{enumerate}
  \item 标变价：标出发生变化的元素化合价
  \item 列变化：列出化合价升降数值
  \item 求最小公倍数：使化合价升降总数相等
  \item 配系数：先配氧化剂、还原剂系数，再配其他
\end{enumerate}

\textbf{常见氧化剂还原性顺序}：
\[
\text{氧化性}：\text{Cl}_2 > \text{Br}_2 > \text{Fe}^{3+} > \text{I}_2 > \text{SO}_2 > \text{S}
\]
\[
\text{还原性}：\text{S}^{2-} > \text{SO}_3^{2-} > \text{I}^- > \text{Fe}^{2+} > \text{Br}^- > \text{Cl}^-
\]

\textbf{常见氧化剂及其产物}：
\[
\begin{array}{c|c|c}
\text{氧化剂} & \text{酸性条件还原产物} & \text{备注} \\ \hline
\text{KMnO}_4 & \text{Mn}^{2+} & \text{紫红色褪去} \\
\text{K}_2\text{Cr}_2\text{O}_7 & \text{Cr}^{3+} & \text{橙色变绿色} \\
\text{HNO}_3(\text{稀}) & \text{NO} & \text{} \\
\text{HNO}_3(\text{浓}) & \text{NO}_2 & \text{} \\
\text{H}_2\text{SO}_4(\text{浓}) & \text{SO}_2 & \text{} \\
\text{H}_2\text{O}_2 & \text{H}_2\text{O} & \text{绿色氧化剂} \\
\text{Fe}^{3+} & \text{Fe}^{2+} & \text{} \\
\end{array}
\]

\subsection{离子反应}

\textbf{离子方程式书写步骤}：
\begin{enumerate}
  \item 写：写出化学方程式
  \item 拆：将可溶性强电解质拆成离子形式
  \item 删：删去两边相同离子
  \item 查：检查电荷守恒、质量守恒
\end{enumerate}

\textbf{需拆的物质}：强酸（$\text{HCl},\text{H}_2\text{SO}_4,\text{HNO}_3$）、强碱（$\text{NaOH},\text{KOH},\text{Ba(OH)}_2$）、可溶性盐

\textbf{不拆的物质}：单质、气体、沉淀、弱电解质（$\text{H}_2\text{O},\text{CH}_3\text{COOH},\text{NH}_3\cdot\text{H}_2\text{O}$）、氧化物

\textbf{常见离子检验}：
\[
\begin{array}{c|c|c}
\text{离子} & \text{试剂} & \text{现象} \\ \hline
\text{Cl}^- & \text{AgNO}_3+\text{HNO}_3 & \text{白色沉淀} \\
\text{SO}_4^{2-} & \text{BaCl}_2+\text{HCl} & \text{白色沉淀} \\
\text{NH}_4^+ & \text{NaOH}加热 & \text{刺激性气体，使湿润红色石蕊试纸变蓝} \\
\text{Fe}^{3+} & \text{KSCN} & \text{溶液变红色} \\
\text{Fe}^{2+} & \text{K}_3[\text{Fe(CN)}_6] & \text{蓝色沉淀} \\
\text{CO}_3^{2-} & \text{稀HCl}+\text{澄清石灰水} & \text{气体使石灰水变浑浊} \\
\end{array}
\]

\subsection{元素周期律与元素周期表}

\textbf{周期表结构}：7 个周期（短周期 1--3，长周期 4--7），16 个族（7 主族 + 7 副族 + Ⅷ族 + 0 族）

\textbf{同周期从左到右}（金属性减弱，非金属性增强）：
\[
\begin{array}{c|c}
\text{性质} & \text{变化趋势} \\ \hline
\text{原子半径} & \text{减小} \\
\text{金属性} & \text{减弱} \\
\text{非金属性} & \text{增强} \\
\text{最高价氧化物对应水化物酸性} & \text{增强} \\
\text{最高价氧化物对应水化物碱性} & \text{减弱} \\
\text{气态氢化物稳定性} & \text{增强} \\
\end{array}
\]

\textbf{同主族从上到下}（金属性增强，非金属性减弱）：
\[
\begin{array}{c|c}
\text{性质} & \text{变化趋势} \\ \hline
\text{原子半径} & \text{增大} \\
\text{金属性} & \text{增强} \\
\text{非金属性} & \text{减弱} \\
\text{最高价氧化物对应水化物酸性} & \text{减弱} \\
\text{最高价氧化物对应水化物碱性} & \text{增强} \\
\text{气态氢化物稳定性} & \text{减弱} \\
\end{array}
\]

\textbf{微粒半径比较规律}：
\begin{itemize}
  \item 同周期：原子半径随原子序数增大而减小
  \item 同主族：原子半径随原子序数增大而增大
  \item 同元素：阴离子半径 $>$ 原子半径 $>$ 阳离子半径
  \item 电子层结构相同的离子：核电荷数越大，半径越小
\end{itemize}

\textbf{易错}：
\begin{itemize}
  \item 比较非金属性强弱看最高价含氧酸酸性，不是看无氧酸酸性
  \item 稀有气体不参与比较（已达稳定结构）
\end{itemize}

\subsection{化学反应与能量}

\textbf{焓变}：
\[
\Delta H = H(\text{生成物}) - H(\text{反应物})
\]
放热反应 $\Delta H < 0$，吸热反应 $\Delta H > 0$

\textbf{热化学方程式书写要点}：
\begin{enumerate}
  \item 标明物质状态（s, l, g, aq）
  \item $\Delta H$ 单位 $\text{kJ/mol}$，与系数对应
  \item 系数加倍，$\Delta H$ 加倍；反应逆向，$\Delta H$ 变号
\end{enumerate}

\textbf{盖斯定律}：
\[
\Delta H = \Delta H_1 + \Delta H_2 + \Delta H_3 + \cdots
\]

\textbf{键能与焓变}：
\[
\Delta H = \sum E(\text{反应物键能}) - \sum E(\text{生成物键能})
\]

\textbf{燃烧热}：$1\ \text{mol}$ 纯物质完全燃烧生成稳定氧化物所放出的热量

\textbf{中和热}：$57.3\ \text{kJ/mol}$（稀强酸与稀强碱反应生成 $1\ \text{mol}\ \text{H}_2\text{O}$）

\textbf{原电池}：
\begin{itemize}
  \item 负极：失电子，发生氧化反应
  \item 正极：得电子，发生还原反应
  \item 电子流向：负极 $\to$ 导线 $\to$ 正极
  \item 电流方向：正极 $\to$ 导线 $\to$ 负极
\end{itemize}

\textbf{电解池}：
\begin{itemize}
  \item 阳极：失电子，发生氧化反应（与电源正极相连）
  \item 阴极：得电子，发生还原反应（与电源负极相连）
  \item 阳离子向阴极移动，阴离子向阳极移动
\end{itemize}

\textbf{金属腐蚀与防护}：
\begin{itemize}
  \item 电化学腐蚀：吸氧腐蚀为主（中性/碱性），析氢腐蚀（酸性）
  \item 防护方法：覆盖保护层、牺牲阳极保护法、外加电流阴极保护法
\end{itemize}

\subsection{化学反应速率与化学平衡}

\textbf{速率公式}：
\[
v = \frac{\Delta c}{\Delta t}\ (\text{mol}/(\text{L}\cdot\text{s}))
\]

\textbf{影响速率的因素}：
\begin{itemize}
  \item 浓度：增大浓度，速率加快
  \item 压强：增大压强（缩小体积），速率加快（涉及气体）
  \item 温度：升高温度，速率加快
  \item 催化剂：正催化剂加快，负催化剂减慢
  \item 其他：固体表面积、浓度梯度等
\end{itemize}

\textbf{勒夏特列原理}：改变影响平衡的一个条件，平衡向减弱这种改变的方向移动。

\textbf{平衡常数}：
\[
aA + bB \rightleftharpoons cC + dD,\qquad K = \frac{[C]^c[D]^d}{[A]^a[B]^b}
\]

$K$ 只与温度有关，$K>10^5$ 反应进行完全。

\textbf{转化率}：
\[
\alpha = \frac{\text{已转化的反应物量}}{\text{反应物起始量}} \times 100\%
\]

\textbf{几种典型平衡}：
\begin{itemize}
  \item $2\text{NO}_2 \rightleftharpoons \text{N}_2\text{O}_4$（压强增大，颜色先深后浅）
  \item $\text{N}_2 + 3\text{H}_2 \rightleftharpoons 2\text{NH}_3$（合成氨，$\Delta H<0$）
\end{itemize}

\textbf{等效平衡}：
\begin{itemize}
  \item 恒温恒容：投料成比例 $\to$ 等效
  \item 恒温恒压：投料成比例 $\to$ 等效
\end{itemize}

\textbf{易错}：
\begin{itemize}
  \item 固体、纯液体的浓度视为常数，不写入 $K$ 表达式
  \item 催化剂同等改变正逆反应速率，不改变平衡
  \item 压强的改变必须引起浓度变化才影响平衡
\end{itemize}

\subsection{电解质溶液}

\textbf{强弱电解质}：
\[
\begin{array}{c|c}
\text{强电解质} & \text{弱电解质} \\ \hline
\text{完全电离} & \text{部分电离} \\
\text{强酸、强碱、大部分盐} & \text{弱酸、弱碱、水} \\
\text{导电性强（同浓度）} & \text{导电性弱（同浓度）} \\
\end{array}
\]

\textbf{电离常数}：
\[
K_a = \frac{[\text{H}^+][\text{A}^-]}{[\text{HA}]},\qquad K_b = \frac{[\text{NH}_4^+][\text{OH}^-]}{[\text{NH}_3\cdot\text{H}_2\text{O}]}
\]

\textbf{水的离子积}：
\[
K_w = [\text{H}^+][\text{OH}^-] = 1\times10^{-14}\ (25^\circ\text{C})
\]

\textbf{pH}：
\[
\text{pH} = -\lg[\text{H}^+],\qquad \text{pOH} = -\lg[\text{OH}^-],\qquad \text{pH}+\text{pOH}=14
\]

\textbf{盐类水解规律}：谁弱谁水解，谁强显谁性
\begin{itemize}
  \item 强酸弱碱盐 $\to$ 酸性（如 $\text{NH}_4\text{Cl}$）
  \item 强碱弱酸盐 $\to$ 碱性（如 $\text{CH}_3\text{COONa}$）
  \item 强酸强碱盐 $\to$ 中性（如 $\text{NaCl}$）
  \item 弱酸弱碱盐 $\to$ 谁强显谁性
\end{itemize}

\textbf{溶度积}：
\[
K_{sp} = [\text{M}^{n+}]^m[\text{A}^{m-}]^n
\]
\begin{itemize}
  \item $Q_c > K_{sp}$：析出沉淀
  \item $Q_c = K_{sp}$：饱和溶液
  \item $Q_c < K_{sp}$：未饱和
\end{itemize}

\textbf{常用滴定指示剂}：
\[
\begin{array}{c|c|c}
\text{指示剂} & \text{变色范围(pH)} & \text{颜色变化} \\ \hline
\text{甲基橙} & 3.1\text{--}4.4 & \text{红 $\to$ 黄} \\
\text{酚酞} & 8.2\text{--}10.0 & \text{无色 $\to$ 红} \\
\text{石蕊} & 5.0\text{--}8.0 & \text{红 $\to$ 蓝} \\
\end{array}
\]

\subsection{非金属元素}

\textbf{氯及其化合物}：
\[
\text{Cl}_2 + 2\text{NaOH} = \text{NaCl} + \text{NaClO} + \text{H}_2\text{O}
\]
\[
\text{Cl}_2 + \text{H}_2\text{O} \rightleftharpoons \text{HCl} + \text{HClO}
\]
\[
\text{MnO}_2 + 4\text{HCl}(\text{浓}) \xrightarrow{\Delta} \text{MnCl}_2 + \text{Cl}_2\uparrow + 2\text{H}_2\text{O}
\]

\textbf{氧族元素（硫）}：
\[
\text{S} + \text{O}_2 \xrightarrow{\text{点燃}} \text{SO}_2
\]
\[
2\text{SO}_2 + \text{O}_2 \xrightarrow[\Delta]{\text{催化剂}} 2\text{SO}_3
\]
\[
\text{Cu} + 2\text{H}_2\text{SO}_4(\text{浓}) \xrightarrow{\Delta} \text{CuSO}_4 + \text{SO}_2\uparrow + 2\text{H}_2\text{O}
\]

\textbf{氮及其化合物}：
\[
\text{N}_2 + 3\text{H}_2 \xrightarrow[\text{高温高压}]{\text{催化剂}} 2\text{NH}_3
\]
\[
4\text{NH}_3 + 5\text{O}_2 \xrightarrow[\Delta]{\text{Pt}} 4\text{NO} + 6\text{H}_2\text{O}
\]
\[
3\text{Cu} + 8\text{HNO}_3(\text{稀}) = 3\text{Cu}(\text{NO}_3)_2 + 2\text{NO}\uparrow + 4\text{H}_2\text{O}
\]
\[
\text{Cu} + 4\text{HNO}_3(\text{浓}) = \text{Cu}(\text{NO}_3)_2 + 2\text{NO}_2\uparrow + 2\text{H}_2\text{O}
\]

\textbf{碳族元素}：
\[
\text{SiO}_2 + 2\text{NaOH} = \text{Na}_2\text{SiO}_3 + \text{H}_2\text{O}
\]
\[
\text{SiO}_2 + 4\text{HF} = \text{SiF}_4\uparrow + 2\text{H}_2\text{O}
\]

\textbf{易错}：
\begin{itemize}
  \item $\text{SO}_2$ 使品红褪色是漂白性（可逆），使 $\text{KMnO}_4$ 褪色是还原性
  \item 浓硝酸与金属反应不产生 $\text{H}_2$
  \item 常温下 $\text{Al}$、$\text{Fe}$ 在浓 $\text{HNO}_3$ 和浓 $\text{H}_2\text{SO}_4$ 中钝化
\end{itemize}

\subsection{金属元素}

\textbf{钠及其化合物}：
\[
2\text{Na} + \text{O}_2 \xrightarrow{\text{点燃}} \text{Na}_2\text{O}_2\ (\text{淡黄色})
\]
\[
2\text{Na}_2\text{O}_2 + 2\text{H}_2\text{O} = 4\text{NaOH} + \text{O}_2\uparrow
\]
\[
2\text{Na}_2\text{O}_2 + 2\text{CO}_2 = 2\text{Na}_2\text{CO}_3 + \text{O}_2
\]

\textbf{铝及其化合物}：
\[
2\text{Al} + 2\text{NaOH} + 2\text{H}_2\text{O} = 2\text{NaAlO}_2 + 3\text{H}_2\uparrow
\]
\[
\text{Al}_2\text{O}_3 + 2\text{NaOH} = 2\text{NaAlO}_2 + \text{H}_2\text{O}
\]
\[
\text{Al(OH)}_3 + \text{NaOH} = \text{NaAlO}_2 + 2\text{H}_2\text{O}
\]

\textbf{铁及其化合物}：
\[
3\text{Fe} + 4\text{H}_2\text{O}(g) \xrightarrow{\Delta} \text{Fe}_3\text{O}_4 + 4\text{H}_2
\]
\[
2\text{Fe}^{3+} + \text{Fe} = 3\text{Fe}^{2+}
\]
\[
2\text{Fe}^{3+} + \text{Cu} = 2\text{Fe}^{2+} + \text{Cu}^{2+}
\]

\textbf{铜及其化合物}：
\[
2\text{Cu} + \text{O}_2 + \text{CO}_2 + \text{H}_2\text{O} = \text{Cu}_2(\text{OH})_2\text{CO}_3\ (\text{铜绿})
\]
\[
\text{CuO} + \text{H}_2 \xrightarrow{\Delta} \text{Cu} + \text{H}_2\text{O}
\]

\textbf{金属活动性顺序}：
\[
\text{K} > \text{Ca} > \text{Na} > \text{Mg} > \text{Al} > \text{Zn} > \text{Fe} > \text{Sn} > \text{Pb} > (\text{H}) > \text{Cu} > \text{Hg} > \text{Ag} > \text{Pt} > \text{Au}
\]

\textbf{金属冶炼方法}：
\[
\begin{array}{c|c|c}
\text{方法} & \text{适用金属} & \text{实例} \\ \hline
\text{电解法} & \text{K}\sim\text{Al} & 2\text{NaCl}(\text{熔融}) \xrightarrow{\text{电解}} 2\text{Na} + \text{Cl}_2 \\
\text{热还原法} & \text{Zn}\sim\text{Cu} & \text{Fe}_2\text{O}_3 + 3\text{CO} \xrightarrow{\Delta} 2\text{Fe} + 3\text{CO}_2 \\
\text{热分解法} & \text{Hg}\sim\text{Ag} & 2\text{HgO} \xrightarrow{\Delta} 2\text{Hg} + \text{O}_2 \\
\text{物理提取} & \text{Pt},\text{Au} & \text{游离态存在} \\
\end{array}
\]

\subsection{有机化学基础}

\textbf{有机反应类型}：
\[
\begin{array}{c|c}
\text{反应类型} & \text{实例} \\ \hline
\text{取代反应} & \text{CH}_4 + \text{Cl}_2 \xrightarrow{\text{光照}} \text{CH}_3\text{Cl} + \text{HCl} \\
\text{加成反应} & \text{CH}_2=\text{CH}_2 + \text{Br}_2 \to \text{CH}_2\text{BrCH}_2\text{Br} \\
\text{消去反应} & \text{C}_2\text{H}_5\text{OH} \xrightarrow{\text{浓}\text{H}_2\text{SO}_4,\ 170^\circ\text{C}} \text{CH}_2=\text{CH}_2\uparrow + \text{H}_2\text{O} \\
\text{氧化反应} & 2\text{CH}_3\text{CH}_2\text{OH} + \text{O}_2 \xrightarrow{\text{Cu},\Delta} 2\text{CH}_3\text{CHO} + 2\text{H}_2\text{O} \\
\text{还原反应} & \text{R}-\text{NO}_2 \xrightarrow{\text{Fe,HCl}} \text{R}-\text{NH}_2 \\
\text{加聚反应} & n\text{CH}_2=\text{CH}_2 \xrightarrow{\text{催化剂}} \text{[CH}_2-\text{CH}_2]_n \\
\text{缩聚反应} & n\text{NH}_2(\text{CH}_2)_5\text{COOH} \to \text{[NH(CH}_2)_5\text{CO]}_n + n\text{H}_2\text{O} \\
\end{array}
\]

\textbf{官能团与特征反应}：
\[
\begin{array}{c|c|c}
\text{官能团} & \text{代表物} & \text{特征反应} \\ \hline
\text{C=C} & \text{乙烯} & \text{加成、加聚、使}\ \text{KMnO}_4\ \text{褪色} \\
\text{—OH（醇）} & \text{乙醇} & \text{酯化、消去、催化氧化} \\
\text{—OH（酚）} & \text{苯酚} & \text{弱酸性、与}\ \text{FeCl}_3\ \text{显紫色} \\
\text{—CHO} & \text{乙醛} & \text{银镜反应、与新制}\ \text{Cu(OH)}_2\ \text{反应} \\
\text{—COOH} & \text{乙酸} & \text{酸性、酯化反应} \\
\text{—COOR} & \text{乙酸乙酯} & \text{水解（酸/碱催化）} \\
\text{—NH}_2 & \text{苯胺} & \text{弱碱性} \\
\text{—X} & \text{溴乙烷} & \text{水解、消去} \\
\end{array}
\]

\textbf{同分异构体书写要点}：
\begin{enumerate}
  \item 碳链异构（主链由长到短）
  \item 位置异构（官能团位置变化）
  \item 官能团异构（不同类别）
\end{enumerate}

\textbf{常见官能团异构}：
\[
\begin{array}{c|c}
\text{分子式} & \text{可能类别} \\ \hline
\text{C}_n\text{H}_{2n} & \text{烯烃/环烷烃} \\
\text{C}_n\text{H}_{2n-2} & \text{炔烃/二烯烃/环烯烃} \\
\text{C}_n\text{H}_{2n+2}\text{O} & \text{醇/醚} \\
\text{C}_n\text{H}_{2n}\text{O} & \text{醛/酮/烯醇/环醚} \\
\text{C}_n\text{H}_{2n}\text{O}_2 & \text{羧酸/酯/羟基醛} \\
\text{C}_n\text{H}_{2n-6}\text{O} & \text{酚/芳香醇/芳香醚} \\
\end{array}
\]

\subsection{化学实验}

\textbf{常用仪器}：
\[
\begin{array}{c|c}
\text{仪器} & \text{用途} \\ \hline
\text{坩埚} & \text{高温灼烧固体} \\
\text{蒸发皿} & \text{蒸发浓缩溶液} \\
\text{蒸馏烧瓶} & \text{蒸馏（需加沸石）} \\
\text{冷凝管} & \text{冷凝蒸气（下口进水）} \\
\text{容量瓶} & \text{配制一定体积的溶液} \\
\text{滴定管} & \text{酸碱中和滴定} \\
\text{分液漏斗} & \text{分液、滴加液体} \\
\text{洗气瓶} & \text{洗气、除杂} \\
\end{array}
\]

\textbf{分离提纯方法}：
\[
\begin{array}{c|c|c}
\text{方法} & \text{适用对象} & \text{实例} \\ \hline
\text{过滤} & \text{固液分离} & \text{沉淀与溶液} \\
\text{蒸发结晶} & \text{可溶性固体} & \text{NaCl 溶液} \\
\text{降温结晶} & \text{溶解度受温度影响大的物质} & \text{KNO}_3 \\
\text{蒸馏} & \text{互溶液体（沸点差异大）} & \text{酒精与水} \\
\text{分液} & \text{互不相溶的液体} & \text{油与水} \\
\text{萃取} & \text{利用溶解度差异} & \text{I}_2\ \text{从水中到}\ \text{CCl}_4 \\
\text{洗气} & \text{气体除杂} & \text{除}\ \text{CO}_2\ \text{中的}\ \text{HCl} \\
\text{盐析} & \text{蛋白质} & \text{加}\ (\text{NH}_4)_2\text{SO}_4 \\
\end{array}
\]

\textbf{常见气体检验}：
\[
\begin{array}{c|c}
\text{气体} & \text{检验方法} \\ \hline
\text{H}_2 & \text{纯氢气安静燃烧，淡蓝色火焰} \\
\text{O}_2 & \text{带火星木条复燃} \\
\text{Cl}_2 & \text{湿润淀粉碘化钾试纸变蓝} \\
\text{NH}_3 & \text{湿润红色石蕊试纸变蓝} \\
\text{CO}_2 & \text{澄清石灰水变浑浊} \\
\text{SO}_2 & \text{品红溶液褪色，加热恢复} \\
\text{NO} & \text{无色气体，遇空气变红棕色} \\
\text{NO}_2 & \text{红棕色气体} \\
\end{array}
\]

\textbf{实验安全与常见事故处理}：
\begin{itemize}
  \item 浓硫酸稀释：酸入水，沿器壁，不断搅拌
  \item 碱液溅到皮肤上：大量水冲洗，涂硼酸溶液
  \item 酒精灯起火：用湿抹布盖灭，不可用水
  \item 温度计水银球位置：蒸馏时在支管口处；测反应液温度时浸入液面下
\end{itemize}

\subsection{物质结构与性质（选修三）}

\textbf{原子结构}：
\[
\text{能级顺序}：1s\ 2s\ 2p\ 3s\ 3p\ 4s\ 3d\ 4p\ 5s\ 4d\ 5p\ 6s\ 4f\ 5d\ 6p\ 7s
\]
\[
\text{洪特规则}：同一亚层电子优先单独占据不同轨道且自旋平行
\]

\textbf{键参数}：
\[
\begin{array}{c|c}
\text{概念} & \text{定义} \\ \hline
\text{键能} & \text{断裂 1 mol 化学键所需的能量} \\
\text{键长} & \text{两原子核间的平均距离} \\
\text{键角} & \text{共价键之间的夹角} \\
\end{array}
\]

\textbf{杂化轨道类型}：
\[
\begin{array}{c|c|c}
\text{杂化方式} & \text{空间构型} & \text{实例} \\ \hline
sp & \text{直线形} & \text{C}_2\text{H}_2,\ \text{CO}_2 \\
sp^2 & \text{平面三角形} & \text{C}_2\text{H}_4,\ \text{SO}_2 \\
sp^3 & \text{四面体形} & \text{CH}_4,\ \text{H}_2\text{O},\ \text{NH}_3 \\
\end{array}
\]

\textbf{分子间作用力}：
\[
\text{范德华力} < \text{氢键} < \text{化学键}
\]

\textbf{晶体类型}：
\[
\begin{array}{c|c|c|c}
\text{晶体类型} & \text{构成微粒} & \text{作用力} & \text{性质} \\ \hline
\text{离子晶体} & \text{阴、阳离子} & \text{离子键} & \text{硬而脆，熔沸点较高} \\
\text{共价晶体} & \text{原子} & \text{共价键} & \text{硬度最大，熔沸点很高} \\
\text{金属晶体} & \text{金属阳离子、自由电子} & \text{金属键} & \text{导电导热，有延展性} \\
\text{分子晶体} & \text{分子} & \text{分子间作用力} & \text{熔沸点低} \\
\end{array}
\]
